% Auto-generated complete chapter from 28C_github_profiles_benchmark.md
\chapter{29\_github\_profiles\_benchmark.md}
\label{complete:root_033}

\begin{sourcemeta}
\textbf{Fuente:} \href{file:///E:/Laboral/28C_github_profiles_benchmark.md}{28C\_github\_profiles\_benchmark.md}\\
\textbf{Seccion:} root\\
\textbf{Estado:} reference\\
\textbf{Rol:} Benchmark o complemento\\
\textbf{Lineas originales:} 737\\
\textbf{Ultima modificacion:} 2026-06-11 17:36\\
\textbf{Deduplicacion conservadora:} 0 bloques repetidos omitidos; 0 caracteres repetidos no reimpresos.
\end{sourcemeta}

\section*{Resumen de orientacion}
This module does \textbf{not} reopen the base labor-market strategy.

\section*{Contenido completo depurado}
\addcontentsline{toc}{section}{Contenido completo depurado}




\subsection{0. Module status}




\begin{mdcode}
Bloque de codigo: text
Module: 29_github_profiles_benchmark
Task: A3_github_profiles_benchmark
Generated: 2026-06-11
Project: Laboral / Alexander Woodcock Salomón
Purpose: Convert GitHub benchmark results into actionable updates for Alexander's public GitHub surface.
\end{mdcode}




This module does \textbf{not} reopen the base labor-market strategy.




Target positioning remains:




\begin{mdcode}
Bloque de codigo: text
Real-Time 3D Developer / Unity Technical Artist
focused on Unity WebGL, CAD-to-realtime optimization, technical visualization, and tooling.
\end{mdcode}




\medskip\hrule\medskip




\subsection{1. Quality audit of A3}




\subsubsection{1.1 Result quality}




A3 is usable.




It collected \textbf{21 GitHub examples}, which satisfies the requested range of \textbf{20–35} profiles/repositories.




The dataset is strongest for:




\begin{mdcode}
Bloque de codigo: text
Unity tooling
Unity technical art utilities
Unity WebGL deployment
XR samples
robotics / digital twin visualization
build automation
LLM-assisted Unity tooling
\end{mdcode}




The dataset is weaker for:




\begin{mdcode}
Bloque de codigo: text
GitHub profile READMEs
Three.js / WebGPU portfolios
general-purpose WebGL visualization
GitHub Pages portfolio examples
\end{mdcode}




This is acceptable because Alexander’s immediate public proof should be Unity/WebGL/tooling-led, not generic WebGL-led.




\medskip\hrule\medskip




\subsection{2. Source-quality classification}




\begin{center}
\scriptsize
\begin{longtable}{>{\raggedright\arraybackslash}p{0.305\linewidth} >{\raggedright\arraybackslash}p{0.305\linewidth} >{\raggedright\arraybackslash}p{0.305\linewidth}}
\toprule
Class & Count & Usefulness \\
\midrule
\endfirsthead
\toprule
Class & Count & Usefulness \\
\midrule
\endhead
Strong & 14 & Directly adaptable \\
Medium & 5 & Partial pattern \\
Weak & 2 & Limited pattern \\
\bottomrule
\end{longtable}
\end{center}




\subsubsection{2.1 Strong examples}




\begin{center}
\scriptsize
\begin{longtable}{>{\raggedright\arraybackslash}p{0.455\linewidth} >{\raggedright\arraybackslash}p{0.455\linewidth}}
\toprule
Label & Why strong \\
\midrule
\endfirsthead
\toprule
Label & Why strong \\
\midrule
\endhead
RAZLUTA-Profile & Strong pinned tooling surface \\
SPLATVFX & Clear experimental graphics README \\
UNITY-BUILDER & Strong tooling docs \\
REACT-UNITY-WEBGL & Strong WebGL package docs \\
UNITY-EDITOR-TOOLBOX & Deep editor-tool docs \\
UNITY-MCP & Strong AI/tooling positioning \\
PY-BUILD-AUTOMATION & Clear Python automation docs \\
UNITY-DATA-TOOLS & Strong analysis-tool README \\
ROS-TCP-CONNECTOR & Strong robotics integration docs \\
XR-INTERACTION-EXAMPLES & Strong XR sample structure \\
TEXTURE-PACKER & Strong technical-art tool docs \\
PACKAGE-MAKER & Strong package-tool docs \\
LIGHTS-LOD & Strong optimization-tool docs \\
ICON-GENERATOR & Strong production-tool docs \\
\bottomrule
\end{longtable}
\end{center}




\subsubsection{2.2 Medium examples}




\begin{center}
\scriptsize
\begin{longtable}{>{\raggedright\arraybackslash}p{0.455\linewidth} >{\raggedright\arraybackslash}p{0.455\linewidth}}
\toprule
Label & Limitation \\
\midrule
\endfirsthead
\toprule
Label & Limitation \\
\midrule
\endhead
VFXGRAPH-TESTBED & Good media, sparse explanation \\
REACHY2-DT & Relevant, but hardware-specific \\
PDT-UNITY & Relevant, academic/research-heavy \\
RUNTIME-UNITY-EDITOR & Strong tool, modding-specific \\
LIGHTS-AUDIT & Useful, narrower scope \\
\bottomrule
\end{longtable}
\end{center}




\subsubsection{2.3 Weak examples}




\begin{center}
\scriptsize
\begin{longtable}{>{\raggedright\arraybackslash}p{0.455\linewidth} >{\raggedright\arraybackslash}p{0.455\linewidth}}
\toprule
Label & Limitation \\
\midrule
\endfirsthead
\toprule
Label & Limitation \\
\midrule
\endhead
VR-PROJECT-TEMPLATE & Template-heavy \\
VR-GAME-JAM-TEMPLATE & README too sparse \\
\bottomrule
\end{longtable}
\end{center}




\medskip\hrule\medskip




\subsection{3. Corrected source inventory}




Some A3 rows contained citation artifacts or uncertain license wording. This inventory normalizes the usable sources.




\begin{center}
\scriptsize
\begin{longtable}{>{\raggedright\arraybackslash}p{0.184\linewidth} >{\raggedright\arraybackslash}p{0.184\linewidth} >{\raggedright\arraybackslash}p{0.184\linewidth} >{\raggedright\arraybackslash}p{0.184\linewidth} >{\raggedright\arraybackslash}p{0.184\linewidth}}
\toprule
Label & URL & Type & Role relevance & Quality \\
\midrule
\endfirsthead
\toprule
Label & URL & Type & Role relevance & Quality \\
\midrule
\endhead
RAZLUTA-Profile & \url{https://github.com/razluta} & Profile & Unity tools / technical art & Strong \\
VFXGRAPH-TESTBED & \url{https://github.com/keijiro/VfxGraphTestbed} & Repo & Unity VFX & Medium \\
SPLATVFX & \url{https://github.com/keijiro/SplatVFX} & Repo & VFX / splats & Strong \\
UNITY-BUILDER & \url{https://github.com/game-ci/unity-builder} & Repo & Build pipeline & Strong \\
REACT-UNITY-WEBGL & \url{https://github.com/jeffreylanters/react-unity-webgl} & Repo & Unity WebGL & Strong \\
REACHY2-DT & \url{https://github.com/pollen-robotics/Reachy2-UnityDigitalTwin} & Repo & XR / digital twin & Medium \\
PDT-UNITY & \url{https://github.com/programming-digital-twins/LabBenchStudios-PDT-Unity} & Repo & Digital twin & Medium \\
RUNTIME-UNITY-EDITOR & \url{https://github.com/ManlyMarco/RuntimeUnityEditor} & Repo & Runtime tooling & Medium \\
UNITY-EDITOR-TOOLBOX & \url{https://github.com/arimger/Unity-Editor-Toolbox} & Repo & Editor tooling & Strong \\
UNITY-MCP & \url{https://github.com/CoplayDev/unity-mcp} & Repo & AI + Unity tooling & Strong \\
PY-BUILD-AUTOMATION & \url{https://github.com/angrysharkstudio/Unity-Python-Build-Automation} & Repo & Python build automation & Strong \\
UNITY-DATA-TOOLS & \url{https://github.com/Unity-Technologies/UnityDataTools} & Repo & Build analysis & Strong \\
ROS-TCP-CONNECTOR & \url{https://github.com/Unity-Technologies/ROS-TCP-Connector} & Repo & Robotics visualization & Strong \\
XR-INTERACTION-EXAMPLES & \url{https://github.com/Unity-Technologies/XR-Interaction-Toolkit-Examples} & Repo & XR samples & Strong \\
VR-PROJECT-TEMPLATE & \url{https://github.com/SamuelAsherRivello/unity-project-template-vr} & Repo & VR template & Weak \\
VR-GAME-JAM-TEMPLATE & \url{https://github.com/ValemVR/VR-Game-Jam-Template} & Repo & VR template & Weak \\
TEXTURE-PACKER & \url{https://github.com/razluta/UnityTextureRgbPacker} & Repo & Technical art utility & Strong \\
PACKAGE-MAKER & \url{https://github.com/razluta/UnityPackageMaker} & Repo & Package tooling & Strong \\
LIGHTS-LOD & \url{https://github.com/razluta/UnityLightsLodSystem} & Repo & Lighting optimization & Strong \\
LIGHTS-AUDIT & \url{https://github.com/razluta/UnityLightsAuditTool} & Repo & Lighting audit & Medium \\
ICON-GENERATOR & \url{https://github.com/razluta/UnityIconGenerationFromModels} & Repo & Asset pipeline utility & Strong \\
\bottomrule
\end{longtable}
\end{center}




\medskip\hrule\medskip




\subsection{4. Benchmark patterns}




\subsubsection{4.1 Strong GitHub surfaces show value immediately}




The strongest repositories do not start with long academic explanations.




They show:




\begin{mdcode}
Bloque de codigo: text
title
one-line value proposition
badges
hero image / GIF
live demo or docs link
quick-start path
\end{mdcode}




For Alexander, the top of the README should answer:




\begin{mdcode}
Bloque de codigo: text
What is this?
What does it demonstrate?
Can I open it?
Can I see screenshots?
Can I understand the technical stack?
\end{mdcode}




\medskip\hrule\medskip




\subsubsection{4.2 Strong READMEs separate user-facing value from implementation}




Good structure:




\begin{mdcode}
Bloque de codigo: text
What it does
Live demo
Key features
Screenshots / GIFs
Technical stack
Architecture
Setup
Validation / metrics
Scope limits
License
\end{mdcode}




Weak structure:




\begin{mdcode}
Bloque de codigo: text
long thesis summary
unfiltered internal notes
unclear scope
no media
no demo path
no license boundary
\end{mdcode}




\medskip\hrule\medskip




\subsubsection{4.3 Strong technical-art tools show before/after or workflow}




For technical art and tooling, the benchmark consistently favors:




\begin{mdcode}
Bloque de codigo: text
GIF of tool in use
input/output example
installation path
manual usage
API usage
limitations
license
\end{mdcode}




Alexander should use this pattern for any future utility repo, such as:




\begin{mdcode}
Bloque de codigo: text
Blender batch export helper
Unity WebGL profiling helper
texture/material preparation script
fastener placement or validation tool
UI Toolkit prototype utility
\end{mdcode}




\medskip\hrule\medskip




\subsubsection{4.4 Strong repositories document limitations explicitly}




Good repos openly state:




\begin{mdcode}
Bloque de codigo: text
experimental status
unsupported scope
hardware requirements
Unity version
known limitations
license constraints
\end{mdcode}




This matches Alexander’s thesis repo well because the current project must avoid overclaiming.




The public repo should keep clear boundaries around:




\begin{mdcode}
Bloque de codigo: text
visual product twin
heuristic thermal reading
academic validation
unsupported metrics
hidden/internal managers
non-final prototype modules
\end{mdcode}




\medskip\hrule\medskip




\subsubsection{4.5 Badges help, but only when restrained}




Useful badges for Alexander:




\begin{mdcode}
Bloque de codigo: text
Unity version
WebGL
C#
Academic thesis
GitHub Pages
License / All rights reserved
\end{mdcode}




Avoid:




\begin{mdcode}
Bloque de codigo: text
badge spam
fake quality badges
unused package-manager badges
npm badges unless a package exists
CI badges unless CI exists
\end{mdcode}




\medskip\hrule\medskip




\subsection{5. Current repository audit}




Repository inspected:




\begin{mdcode}
Bloque de codigo: text
delarge95/WebGL-Thesis-Proposal
Visibility: public
Default branch: master
Size: ~41 MB
README status: maintained
\end{mdcode}




\subsubsection{5.1 Current strengths}




The current README already has several benchmark-aligned features:




\begin{mdcode}
Bloque de codigo: text
clear project title
one-line product description
Unity / WebGL / C# / academic badges
public demo link
canonical repository link
report and manual links
public scope disclaimer
canonical folder structure
Unity setup
runtime systems table
validation package explanation
public repository boundary
academic scope
\end{mdcode}




This is better than many public Unity academic repos.




\subsubsection{5.2 Current gaps}




The current README is technically solid but still reads more like a thesis archive than a recruiter-facing GitHub proof page.




Main gaps:




\begin{mdcode}
Bloque de codigo: text
no hero screenshot or GIF near the top
no compact “what this demonstrates” section
no visual feature gallery
no short case-study link block
no technical proof bullets above the fold
no GitHub topics listed in the README
no separate portfolio-facing repo profile
\end{mdcode}




These are not fatal. They are presentation gaps, not structural failures.




\medskip\hrule\medskip




\subsection{6. Recommended updates for WebGL-Thesis-Proposal}




\subsubsection{6.1 Add hero media after badges}




Insert after the badge block:




\begin{mdcode}
Bloque de codigo: md
![TwinSight X500 WebGL demo](docs/media/twinsight-hero.gif)

Live demo: https://delarge95.github.io/WebGL-Thesis-Proposal/
\end{mdcode}




Acceptable alternatives:




\begin{mdcode}
Bloque de codigo: text
animated GIF of selection flow
short MP4 converted to GIF
static 16:9 screenshot if GIF is too heavy
\end{mdcode}




Required visual states:




\begin{mdcode}
Bloque de codigo: text
full drone view
part selection
bottom sheet
exploded view
cross-section
Studio visual mode
\end{mdcode}




\medskip\hrule\medskip




\subsubsection{6.2 Add a compact proof section}




Recommended section:




\begin{mdcode}
Bloque de codigo: md
## What this demonstrates

- Unity WebGL deployment of an interactive technical visualization prototype.
- CAD-to-realtime asset preparation for a complex drone assembly.
- Mobile-first technical UI built around part inspection and spatial understanding.
- Runtime systems for selection, visibility, exploded view, cross-section and visual modes.
- Academic validation package with technical profiling, SUS, NASA-TLX Raw and Think-Aloud protocols.
\end{mdcode}




This should appear before the long thesis summary.




\medskip\hrule\medskip




\subsubsection{6.3 Add visual feature gallery}




Recommended structure:




\begin{mdcode}
Bloque de codigo: md
## Feature gallery

| Feature | Preview | Purpose |
|---|---|---|
| Selection | image | Identify parts |
| Inspect | image | Isolate components |
| Analyze | image | Explode / cut assembly |
| Studio | image | Compare visual readings |
| Technical sheets | image | Read specifications |
\end{mdcode}




Keep cells short. Use 5 rows maximum.




\medskip\hrule\medskip




\subsubsection{6.4 Add repository topics on GitHub}




Recommended GitHub topics:




\begin{mdcode}
Bloque de codigo: text
unity
webgl
unity-webgl
technical-visualization
technical-art
real-time-3d
cad-to-realtime
drone
ui-toolkit
academic-thesis
\end{mdcode}




Optional later:




\begin{mdcode}
Bloque de codigo: text
webassembly
interactive-3d
visual-product-twin
\end{mdcode}




\medskip\hrule\medskip




\subsubsection{6.5 Keep the scope disclaimer}




Do \textbf{not} remove the current visual product twin disclaimer.




It protects against overclaiming. It is also aligned with strong open-source documentation practice.




Keep explicit wording around:




\begin{mdcode}
Bloque de codigo: text
not live telemetry
not calibrated thermal analysis
metrics only when validated
hidden/internal systems not final public features
\end{mdcode}




\medskip\hrule\medskip




\subsubsection{6.6 Do not over-convert the thesis repo into a personal portfolio}




The repository should remain the canonical academic/project repo.




Alexander’s personal GitHub surface should instead be improved through:




\begin{mdcode}
Bloque de codigo: text
profile README
pinned repositories
separate portfolio landing page
TwinSight case-study page
small tooling repos
\end{mdcode}




Do not mix personal biography, job-seeking language and thesis archive text too aggressively inside the thesis repo.




\medskip\hrule\medskip




\subsection{7. Recommended GitHub profile README}




Create or update:




\begin{mdcode}
Bloque de codigo: text
delarge95/delarge95/README.md
\end{mdcode}




Recommended structure:




\begin{mdcode}
Bloque de codigo: md
# Alexander Woodcock Salomón

Real-Time 3D Developer / Unity Technical Artist focused on Unity WebGL, CAD-to-realtime optimization and technical visualization.

## Selected work

- TwinSight X500 — Unity WebGL technical visualization prototype for Holybro X500 V2 drone assembly.
- CAD-to-realtime pipeline — Blender cleanup, UVs, baking, optimization and Unity integration.
- Technical UI systems — selection, bottom sheet, exploded view, cross-section and visual modes.

## Stack

Unity · C# · WebGL · URP · UI Toolkit · Blender · CAD cleanup · Git/GitHub · Python tooling

## Links

- Portfolio: [pending]
- TwinSight demo: https://delarge95.github.io/WebGL-Thesis-Proposal/
- TwinSight repo: https://github.com/delarge95/WebGL-Thesis-Proposal
- LinkedIn: [pending]
\end{mdcode}




Avoid:




\begin{mdcode}
Bloque de codigo: text
student-only framing
long personal history
all tools ever used
unverified metrics
AI-generated badges without function
\end{mdcode}




\medskip\hrule\medskip




\subsection{8. Recommended pinned repositories}




Target pinned set:




\begin{center}
\scriptsize
\begin{longtable}{>{\raggedright\arraybackslash}p{0.305\linewidth} >{\raggedright\arraybackslash}p{0.305\linewidth} >{\raggedright\arraybackslash}p{0.305\linewidth}}
\toprule
Slot & Repo & Purpose \\
\midrule
\endfirsthead
\toprule
Slot & Repo & Purpose \\
\midrule
\endhead
1 & WebGL-Thesis-Proposal & Main proof \\
2 & Portfolio site & Recruiter entry \\
3 & Unity tool repo & Tooling proof \\
4 & Blender/Python automation & Pipeline proof \\
5 & WebGL demo experiment & Deployment proof \\
6 & Shader/visualization study & Technical art proof \\
\bottomrule
\end{longtable}
\end{center}




Current available repo:




\begin{mdcode}
Bloque de codigo: text
WebGL-Thesis-Proposal
\end{mdcode}




Missing repos to create later:




\begin{mdcode}
Bloque de codigo: text
portfolio site
small Unity tool
small Python automation tool
small shader/VFX or WebGL visualization demo
\end{mdcode}




\medskip\hrule\medskip




\subsection{9. Recommended README template for future repos}




Use this template for any future tooling or demo repository:




\begin{mdcode}
Bloque de codigo: md
# Project Name

One-line description of the problem solved.

![demo](docs/media/demo.gif)

## What it does

- Feature 1
- Feature 2
- Feature 3

## Why it matters

Short technical rationale.

## Tech stack

Unity · C# · WebGL · Blender · Python

## Quick start

1. Requirement
2. Installation
3. Run

## Workflow

Input -> Process -> Output

## Screenshots

| Step | Preview |
|---|---|
| Input | image |
| Output | image |

## Limitations

- Known limitation 1
- Known limitation 2

## License

State the license or rights boundary clearly.
\end{mdcode}




\medskip\hrule\medskip




\subsection{10. Concrete updates for existing strategy files}




\subsubsection{10.1 Update GitHub strategy module}




If a module similar to \texttt{18\_github\_profile\_and\_repo\_strategy.md} exists, update it with:




\begin{mdcode}
Bloque de codigo: text
- GitHub profile README is required.
- TwinSight repo remains academic/project canonical.
- Add hero media to TwinSight README.
- Add “What this demonstrates” section.
- Add feature gallery.
- Add GitHub topics.
- Create at least 2 small utility repos after thesis delivery.
\end{mdcode}




\subsubsection{10.2 Update CV module}




In \texttt{17\_cv\_base\_and\_role\_variants.md}, add GitHub proof links:




\begin{mdcode}
Bloque de codigo: text
TwinSight demo
TwinSight repository
TwinSight case study
GitHub profile
\end{mdcode}




Add a line under the main project:




\begin{mdcode}
Bloque de codigo: text
Public GitHub repository includes WebGL build, final report package, manuals, validation instruments and runtime architecture documentation.
\end{mdcode}




\subsubsection{10.3 Update portfolio module}




The portfolio should not duplicate the GitHub README.




Portfolio case study should show:




\begin{mdcode}
Bloque de codigo: text
problem
process
screenshots
technical decisions
final demo
results / validation
\end{mdcode}




GitHub README should show:




\begin{mdcode}
Bloque de codigo: text
how to inspect the repo
how to run / understand the demo
technical stack
scope limits
\end{mdcode}




\medskip\hrule\medskip




\subsection{11. Priority order after A3}




Do not start job applications yet.




Recommended order:




\begin{mdcode}
Bloque de codigo: text
1. Convert A3 into GitHub public-surface checklist.
2. Add missing visual/media assets to TwinSight README.
3. Create GitHub profile README.
4. Benchmark ArtStation / portfolio / case-study pages.
5. Build portfolio case study.
6. Then move into live jobs and companies.
\end{mdcode}




\medskip\hrule\medskip




\subsection{12. Next task}




Next task:




\begin{mdcode}
Bloque de codigo: text
A4_portfolio_artstation_case_study_benchmark
\end{mdcode}




Mode:




\begin{mdcode}
Bloque de codigo: text
Agent: Sí
Deep Research: No
\end{mdcode}




Purpose:




\begin{mdcode}
Bloque de codigo: text
Collect public portfolio, ArtStation, personal-site and case-study examples relevant to real-time 3D, Unity technical art, technical visualization, XR/simulation and tool/pipeline development.
\end{mdcode}




\medskip\hrule\medskip




\subsection{13. Prompt for A4}




\begin{mdcode}
Bloque de codigo: text
Execute A4_portfolio_artstation_case_study_benchmark.

Objective:
Collect 20–30 public portfolio or case-study examples relevant to Alexander’s positioning:
Real-Time 3D Developer, Unity Technical Artist, Unity WebGL Developer, Technical Visualization Developer, XR Developer, Simulation Developer, Tools/Pipeline Developer, Python automation/tooling developer.

Prioritize:
- personal portfolio websites;
- ArtStation project pages;
- Behance case studies if technical enough;
- GitHub Pages portfolios;
- itch.io project pages with technical breakdowns;
- Unity/WebGL live demo pages;
- technical art breakdowns;
- XR/simulation/visualization case studies;
- portfolio pages with strong project cards and media.

Output structured tables only.

Fields:
- Source label
- URL
- Source type
- Role relevance
- Country/region if visible
- Portfolio structure
- Hero section
- Project card structure
- Case study depth
- Media used
- Demo links
- GitHub/source links
- Technical breakdown quality
- Metrics/results shown
- Contact/CTA handling
- What Alexander can adapt
- What Alexander should not copy
- Confidence
- Date checked

Rules:
- Public sources only.
- Do not invent unavailable data.
- If unavailable, write “not found”.
- Keep cells short.
- Include source URLs.
- Tables only.
- No narrative strategy.
\end{mdcode}



